\section{Sur-détermination cohomologique du résidu $1/8$}
\label{sec:surdetermination_un_huitieme}

La densité semi-classique PT diffère de la densité Riemann--von Mangoldt
par une constante $1/8$ (\S\ref{sec:densite_semi_classique}).
\citeauthor{BerryKeating1999} \citeyear{BerryKeating1999} identifient
cette constante comme la phase Maslov fine des coins de la double
barrière, sans interprétation dynamique. PT fournit une \emph{triple
identification cohomologique} convergente, sur-déterminée par la
cohérence des trois lectures.

\subsection{K2 — Identité algébrique : $1/8 = \sper^{2}/2 = \lambdaH$}
\label{sec:K2}

L'auto-couplage scalaire du Higgs dans le secteur PT est dérivé dans la
monographie (ch.~15, profondeur $d = 3$) \cite{PT_Foundations}~:
\begin{equation}
\lambdaH \;=\; \dfrac{\sper^{2}}{2}.
\label{eq:lambdaH}
\end{equation}
À $\sper = 1/2$, $\lambdaH = 1/8$. Cette identité est structurellement
liée à la bifurcation $\qplus/\qminus$ à deux branches~: $\lambdaH$ est
le carré du paramètre de bifurcation divisé par le nombre de branches
scalaires.

\paragraph{Identifications algébriques associées.}
Deux autres identifications algébriques coïncident~:
\[
\dfrac{1}{8} \;=\; \bigg(\dfrac{1}{2}\bigg)^{3}
\quad \text{(cube de l'unit\'e scalaire, \cite{PT_M5})},
\]
\[
\dfrac{1}{8} \;=\; \dfrac{1}{2 N_{\mathrm{gen}} + 2}
\quad \text{(\`a $N_{\mathrm{gen}} = 3$, exact)}.
\]

\paragraph{Candidats naïfs rejetés.}
Plusieurs candidats sont rejetés par incompatibilité numérique~:
$1/N_{\mathrm{Weyl}} = 1/15 \ne 1/8$,
$\chi(\Sigma_{\mathrm{pers}})/Z = -26/Z \ne 1/8$,
$\gammathree \gammafive \gammaseven \approx 0{,}335 \ne 1/8$,
$\alpha_{\mathrm{bare}} \approx 1/136 \ne 1/8$. Le $1/8$ n'est pas un
invariant topologique de la surface de persistance, mais un invariant
scalaire de la bifurcation $\qplus/\qminus$.

\subsection{K3 — Identité cohomologique : $1/8 = c_1/N_{\mathrm{corners}}$}
\label{sec:K3}

La structure cohomologique pertinente est construite dans le corpus PT.
La variété de Fisher doublée
\begin{equation}
\mathcal{M}_m^{\mathrm{db}} \;=\; \mathcal{M}_m \times \mathcal{M}_m, \qquad
g(p,q) \;=\; g_{\mathrm{Fisher}}(p) \,\oplus\, g_{\mathrm{Fisher}}(p),
\label{eq:Mdoub}
\end{equation}
(le \emph{même} point base $p$ dans les deux blocs, construction
canonique de Dombrowski 1962 sur le fibré tangent d'une variété
hessienne) munie de la structure complexe $J_{\mathrm{PT}}(u, v) = (-v, u)$
est une variété kählérienne (théorème~\texttt{kahler\_fisher},
\cite{PT_M5}). Son potentiel kählérien est l'entropie de Shannon
$K(p) = \sum_k p_k \log p_k$ (potentiel hessien canonique de Fisher,
Amari-Nagaoka 2000). Les projecteurs chiraux
\begin{equation}
P_{\pm} \;=\; \tfrac{1}{2}(I \mp i J_{\mathrm{PT}})
\label{eq:Pchir}
\end{equation}
décomposent $H^{*}(BG_m; \Cplx)$ en sous-espaces propres $\pm i$ et sont
identifiés au projecteur spinoriel $\qplus/\qminus$ du crible
(\cite{PT_M5}, \texttt{def:chiral\_projectors}).

\begin{theoreme}[Première classe de Chern du fibré spinoriel PT, \texttt{berry\_3pi}]
\label{thm:berry_3pi}
La première classe de Chern du fibré spinoriel sur
$\mathcal{M}_m^{\mathrm{db}}$ vaut
\begin{equation}
c_1 \;=\; \sper \;=\; \dfrac{1}{2}.
\label{eq:c1}
\end{equation}
L'holonomie du revêtement $N_{\mathrm{gen}}$-fois est
$\exp(2\pi i\,c_1\,N_{\mathrm{gen}}) = \exp(3\pi i)$ (origine du
théorème \texttt{berry\_3pi} de la monographie, app.~P, \S~C7).
\end{theoreme}

Le sous-énoncé algébrique $c_1 = 1/2$ est \textbf{[Thm]} inconditionnel ;
la rigueur APS stricte du calcul $1/8 = c_1/N_{\mathrm{corners}}$ sur la
variété à coin reste partiellement ouverte (\S\ref{sec:limites_ouvertures}),
statut \textbf{[Thm partiel]}.

\begin{proposition}[Calcul cohomologique du résidu $1/8$]
\label{prop:c1_sur_N}
Dans l'espace de phases canonique $(u, p_u)$, la double barrière BK est
un quart de plan symplectique
$Q_\star = \{u \ge u_\star,\, p_u \ge p_\star,\, u_\star\,p_\star = 2\pi\}$.
Le bord admet deux arêtes se rejoignant au coin. Avec la symétrisation
de Weyl et les deux branches miroirs identifiées par l'antipériodicité
$T_3$ (\S\ref{sec:bord_antiperiodique}), on obtient
\begin{equation}
N_{\mathrm{corners}} \;=\; 4
\label{eq:Ncoins}
\end{equation}
coins distincts sur la cellule complète. La phase de Berry contribuée par
chaque coin, mesurée sur le quart de tour reliant les deux arêtes
adjacentes, est
\begin{equation}
\gamma_{\mathrm{corner}} \;=\; -\dfrac{2\pi\,c_1}{N_{\mathrm{corners}}} \;=\; -\dfrac{\pi}{4},
\label{eq:gamma_corner}
\end{equation}
et la densité de Bohr--Sommerfeld associée est
\begin{equation}
\boxed{\;N_{\mathrm{corner}} \;=\; \dfrac{\gamma_{\mathrm{corner}}}{2\pi} \;=\; -\dfrac{c_1}{N_{\mathrm{corners}}} \;=\; -\dfrac{1}{8}.\;}
\label{eq:c1_sur_N}
\end{equation}
\end{proposition}

\begin{proof}[Esquisse]
Le quart de plan symplectique admet 2 arêtes ($u = u_\star$ et
$p_u = p_\star$) qui se rencontrent au coin
$(u_\star, p_\star) = (\sqrt{2\pi}, \sqrt{2\pi})$. La symétrisation de
Weyl (involution $(u, p_u) \leftrightarrow (p_u, u)$) crée le secteur
miroir, doublant les arêtes et créant deux coins supplémentaires.
L'antipériodicité $\theta = \pi$ (proposition~\ref{prop:antiperiodicite})
identifie deux paires de coins en deux, donnant $N_{\mathrm{corners}} = 4$
au total (sans l'identification, on aurait $N = 8$ et
$c_1/N = 1/16$, contredisant le résidu observé~\eqref{eq:residu_un_huitieme}).
La phase de Berry par coin est calculée par transport parallèle sur le
quart de tour interne~: en variables polaires locales centrées au coin,
$\gamma_{\mathrm{corner}} = \oint A = -2\pi\,c_1/N_{\mathrm{corners}} = -\pi/4$
avec $c_1 = 1/2$. Le passage à la densité Bohr--Sommerfeld divise par
$2\pi$~: $N_{\mathrm{corner}} = -1/8$. Vérification numérique~: précision
machine $1{,}68 \cdot 10^{-14}$ via le script \texttt{pt\_k3\_cohomological\_maslov.py}.
\end{proof}

L'identification $1/8 = c_1/N_{\mathrm{corners}}$ donne donc le $-\pi/8$
par coin de Berry--Keating, mais \emph{dérivé} depuis la structure
cohomologique PT, et non comme artefact géométrique opaque.

\subsection{Phase de Berry par coin : $1/8 = \sper/4$}
\label{sec:phase_berry}

À l'équateur de la sphère de Bloch ($\theta = \pi/2$), la connexion de
Berry d'un spin-$\sper$ est constante~:
$A_\varphi = -\sin^{2}(\pi/4) = -1/2 = -\sper$ par unité d'angle
azimutal. Sur un cycle équatorial complet, la phase Berry est
$\gamma_1 = -\pi = -2\pi\,\sper$. Divisée par les quatre
coins~\eqref{eq:Ncoins}~:
\begin{equation}
\dfrac{\gamma_1}{N_{\mathrm{corners}}} \;=\; -\dfrac{\pi}{4} \;=\; -2\pi \cdot \dfrac{\sper}{4},
\label{eq:s_sur_4}
\end{equation}
soit $-1/8$ par coin en compte d'aire. L'identité algébrique
correspondante est
\begin{equation}
\dfrac{1}{8} \;=\; \dfrac{\sper}{4}.
\label{eq:un_huitieme_sur_4}
\end{equation}

\subsection{Sur-détermination par cohérence}
\label{sec:surdetermination_coherence}

Les trois identifications convergent vers le même nombre à $\sper = 1/2$~:
\begin{equation}
\boxed{\;
\dfrac{1}{8} \;=\; \underbrace{\dfrac{\sper^{2}}{2}}_{\text{K2, Higgs}} \;=\; \underbrace{\dfrac{c_1}{N_{\mathrm{corners}}}}_{\text{K3, Chern}} \;=\; \underbrace{\dfrac{\sper}{4}}_{\text{Berry par coin}}.\;}
\label{eq:triple_identification}
\end{equation}
La cohérence algébrique entre K2 et la phase de Berry par coin se réduit
à
\begin{equation}
\dfrac{\sper^{2}}{2} \;=\; \dfrac{\sper}{4} \quad\Longleftrightarrow\quad \sper\,(2\sper - 1) \;=\; 0.
\label{eq:coherence_K2_Berry}
\end{equation}
Les seules solutions sont $\sper = 0$ (trivial, pas de paramètre actif,
exclu par $\Tone$) et $\sper = 1/2$ (théorème $\Tone$). \textbf{La
cohérence des trois lectures cohomologiques sur-détermine $\sper = 1/2$
modulo l'axiome $\Tone$ supposé} ; elle ne dérive pas $\sper = 1/2$
indépendamment de $\Tone$. C'est une signature structurelle de $\Tone$
par la zone HP-PT du corpus, non un substitut à $\Tone$.

\subsection{Lecture en termes de théorème d'indice}
\label{sec:indice_APS}

L'assemblage final s'interprète comme un théorème d'indice
Atiyah--Patodi--Singer formel \cite{APS1975} sur la variété à coin
$\square$~:
\begin{equation}
\mathrm{ind}\!\big(D_+|_\square\big) \;=\; \int_{\square} \mathrm{ch}(E)\,\mathrm{Td}(\square) \,+\, \dfrac{1}{2}\,\big(\eta_{\partial} + \dim \ker D_\partial\big) \,+\, \sum_{\mathrm{coins}} \dfrac{\delta_{\mathrm{coin}}}{2}.
\label{eq:APS}
\end{equation}
Avec $c_1(E) = 1/2$ et $\delta_{\mathrm{coin}} = \pi/2$ pour un coin
droit, la contribution par coin est $\pi/2 / (2 \cdot 2\pi) = 1/8$, en
accord avec le calcul de Berry direct. L'identification équivalente via
Chern--Simons~:
$\exp(2\pi i\,c_1/N_{\mathrm{corners}}) = \exp(i\pi/4)$, dont l'argument
$\pi/4$ donne en densité $1/8$ par coin.

Une rigueur formelle au sens APS strict (au-delà du calcul de l'aire et
de la vérification numérique) reste à développer
(\S\ref{sec:limites_ouvertures}). Le mécanisme cohomologique est
néanmoins explicitement identifié.

